\begin{document}
\maketitle

\section{\textbf{Appointments}}
\cventry{2019- present}{Postdoctoral scholar, Sali Lab, Department of Bioengineering and Therapeutic Sciences}{University of California San Francisco}{San Francisco}{CA}{}

\section{\textbf{Education}}
\cventry{2013- 2018}{Ph.D, Chemical Engineering}{University of California Santa Barbara}{Santa Barbara}{CA}{}

\cventry{2014- 2016}{Graduate emphasis in Computational Science and Engineering}{University of California Santa Barbara}{Santa Barbara}{CA}{(Specialization in numerical methods for differential equations, and parallel computing)}

\cventry{2008- 2013}{Masters and Bachelors of Technology (Hons.) integrated dual degree, Chemical Engineering}{Indian Institute of Technology Kharagpur}{Kharagpur}{WB, India}{}

\section{\textbf{Research Interests}}
\cvline{}{Protein structure modeling using chemical crosslinking data, protein-protein docking, probabilistic machine learning, representation learning for protein graphs, coarse-graining of moecular dynamics simulations}

\section{\textbf{Skills}}
\cvline{Softwares}{MATLAB, Mathematica, LAMMPS, GROMACS, AMBER, IMP, VMD, PyMol, UCSF Chimera}
\cvline{Frameworks}{PyTorch, PyMC3}
\cvline{Languages}{Python (\textasciitilde 10000 lines), Fortran-90 (\textasciitilde 1000 lines), C++ (\textasciitilde 1800 lines), Bash, TCL}
\cvline{Parallel-platforms}{MPI, OpenMP}

\section{\textbf{Research Experience}}
\cvline{\textbf{UCSF}}{Postdoc advisor: Andrej Sali}
\cvlistitem{\textbf{Protein structure modeling using chemical crosslinks:} \newline I develop computational protein structure models by Integrating homology models with mass-spectrometry-chemical-crosslinking data. Beyond the traditional use of chemical crosslink information as distance restraints in computational modeling, I employ crosslink data in new and innovative ways to learn domain representation and perform fragment assembly in protein complexes. I also use deep learning techniques to develop computationally fast yet thermodynamically accurate atomistic models of crosslinking peptides. Application of these methods include chromosomal maintenance (Smc5/6) and repair (DNA helicase) proteins.}

\vspace{0.1in}
\cvlistitem{\textbf{Whole cell models:} \newline As a part of the Pancreatic Beta-Cell Consortium (collaborating across USC, UCLA, UCSF and SanghaiTech), I use Bayesian networks to combine disparate models of the cell such as continuum scale pharmacokinetic models, particle scale Brownian Dynamics models and statistical network models of enzymatic pathways, to build an overarching \emph{metamodel} of the cell.}

\newpage

\cvline{\textbf{UCSB}}{PhD advisor: M. Scott Shell}
\cvlistitem{\textbf{Thesis: Next-Generation Coarse-Grained Models for Molecular Dynamics Simulations of Fluid Phase Equilibria and Protein Biophysics Using the Relative Entropy} \newline Introduced the local density potential for structurally and thermodynamically accurate coarse-grained molecular dynamics (MD) simulations of liquid-liquid phase separation, which are relatively rare in the chemical engineering literature. \newline \newline Also developed protein backbone models for (a) template-free folding of 200+ residue coarse-grained protein domains using MD simulations and, (b) studying self-assembly in amyloidogenic peptides (as seen in neurodegenerative diseases like Alzheimer's, ALS, etc).}
\vspace{0.1in}

\cvline{\textbf{IIT Kharagpur}}{Undergraduate thesis advisor: Saikat Chakraborty}
\cvlistitem{Reactive-diffusive modeling of human respiration and its application in characterizing the patho-physiology of methemoglobin induced anemia. (Awarded best thesis)}

\section{\textbf{Computational experience}}
\cvlistitem{\textbf{Deep generative models} - Experienced in parameterizing generative models, specifically variational autoencoders, using PyTorch.}
\cvlistitem{\textbf{Bayesian inference and Bayesian networks} - Experienced in building and sampling Bayesian networks for model selection, using PyMC3.}
\cvlistitem{\textbf{Analysis of high dimensional data} - Experienced in analyzing massively high-dimensional molecular dynamics trajectories to extract physically motivated low-dimensional projections and constructing meaningful visual representations.}
\cvlistitem{\textbf{Parallel computing} - Parallelized runtime-critical parts of (PhD) group software using MPI}
\cvlistitem{\textbf{Open-source contribution} - Local density potential and replica exchange tool for LAMMPS.}
\cvlistitem{\textbf{Linux cluster management} - System administrator for (PhD) group Rocks 6.2 cluster. Installed Rocks and wrote post-install scripts to automate safe handling of the network-attached-storage and developed protocols to automate resource-sharing between lab members.}

\section{\textbf{Publications}}
\cvlistitem{B.Raveh, L.Sun, K.L. White, T. Sanyal, J. Tempkin, D. Zheng, K.B. Pilla, J. Singla, C,Wang, J. Zhao, A. Li, N.A. Graham, C. Kesselman, R.C. Stevens and A. Sali, \textit{Bayesian metamodeling of complex biological systems across varying representations}, PNAS (2021), in review}
\cvlistitem{Y. Yu, Z. Ser, T. Sanyal, K. Choi, B. Wan, A. Sali, A. Kentsis, D.J. Patel and X. Zhao, \textit{Integrative analysis reveals unique features of the Smc5/6 complex}, PNAS (202), in review, preprint available on bioRxiv (https://doi.org/10.1101/2020.12.31.424863)}
\cvlistitem{T. Sanyal, J. Mittal and M.S. Shell,\textit{ A hybrid, bottom up, structurally-accurate, G$\bar{o}$-like coarse-grained protein model}. J. of Chem. Phys. (2019) 151, 044111 (\textit{Editor's pick article, 2019})}
\vspace{0.1in}
\cvlistitem{D. Rosenberger, T.Sanyal, M.S. Shell and N.F.A van der Vegt, \textit{Transferability of local density-assisted implicit solvation models for homogeneous fluid mixtures.} J. Chem. Theory Comput. (2019) 15, 2881-2895}
\vspace{0.1in}
\cvlistitem{M.P. Howard, W.F. Reinhart, T. Sanyal, M.S. Shell, A. Nikoubashman and A.Z. Panagiotopoulos, \textit{Evaporation-induced assembly of colloidal crystals.} J. Chem. Phys. (2018) 149(\textbf{9})}
\vspace{0.1in}
\cvlistitem{T. Sanyal and M.S. Shell, \textit{Transferable coarse-grained models of liquid-liquid equilibrium using local density potentials optimized with the relative entropy.} J. Phys. Chem. B (2018) 122 (\textbf{21}), 5678-5693}
\vspace{0.1in}
\cvlistitem{T. Sanyal and M.S. Shell, \textit{Coarse-grained models using local-density potentials optimized with the relative entropy: Application to implicit solvation.} J. Chem. Phys. (2016) \textbf{145}, 034109}
\vspace{0.1in}
\cvlistitem{T. Sanyal and S. Chakraborty, \textit{Multiscale analysis of hypoxemia in methemoglobin-anemia.} Math. Biosci. (2013) \textbf{241}, 167-180}
\vspace{0.1in}
\cvlistitem{T. Sanyal and S. Chakraborty, \textit{Multiscale analysis of simultaneous uptake of two interacting gases in the human lungs and its application to methemoglobin anemia.} Comput. Chem. Engg. (2013) \textbf{59}, 226-242}

\section{\textbf{Conference Presentations}}
\cvlistitem{Berkeley Mini Stat Mech Meeting; Berkeley, CA (Jan. 2019) \textbf{Poster}: \textit{ Coarse-grained models for protein folding and self-assembly with the relative entropy\\}}

\cvlistitem{10$^{\mathrm{th}}$ Annual Amgen Clorox Graduate Student Symposium; UC Santa Barbara, CA (Oct. 2017) \textbf{Talk}: \textit{A new approach to computationally fast coarse-grained models of hydrophobic interactions in fluids}\\}

\cvlistitem{AIChE Annual Meeting; San Francisco, CA (Nov. 2016) \textbf{Talk}: \textit{Improved coarse-grained models through local density potentials, optimized with the relative entropy}\\}

\cvlistitem{Berkeley Mini Stat Mech Meeting; Berkeley, CA (Jan. 2015) \textbf{Poster}: \textit{Building coarse grained models of solvation using local density dependent interactions with the relative entropy\\}}

\cvlistitem{Southern California Simulations in Science Conference; UC Santa Barbara, CA (Oct. 2015) \textbf{Poster}: \textit{Building coarse grained models of solvation using local density dependent interactions with the relative entropy}\\}

\cvlistitem{8$^{\mathrm{th}}$Amgen-Clorox Graduate Student Symposium; UC anta Barbara, CA (Sept. 2015) \textbf{Poster}: \textit{Robust models of coarse-grained interactions using multi-body potentials with the relative entropy}\\}

\cvlistitem{AIChE Annual Meeting; San Francisco, CA (Nov. 2013) \textbf{Poster}: \textit{Stability of mixing limited patterns in isothermal homogenous autocatalytic reactions}}

\section{\textbf{Teaching Assistance}}
\cvline{}{Undergraduate Thermodynamics (Fall 2018), Undergraduate Reaction Engineering (Spring 2017 ), Graduate Thermodynamics (Fall 2015), Introduction to Chemical Engineering(Fall 2014}

\section{\textbf{Undergraduates mentored}}
\cvlistitem{Jun Lu, UCSB, 2018-2019, Project: \textit{Effects of polymer conjugation on tethered peptides}}

\section{\textbf{Relevant Courses}}
\cvline{}{Thermodynamics, Equilibrium and Non-equilibrium statistical mechanics, Finite-Difference and Finite-Element methods for PDEs, Parallel computing}

\section{\textbf{Awards}}
\cvline{2012}{\textbf{Best bachelors thesis} IIT Kharagpur}
\cvline{2009}{\textbf{Technopreneur Promotion Programme, Ministry of Science and Technology, Govt. of India} - seed money for prototype scale-up of novel bio-reactors for algal bio-diesel generation}
\cvline{2009}{\textbf{Jagdis Bose National Science Talent Search (JBNSTS) foundation, WB, India} - fellowship}

\section{\textbf{Outreach Activities}}
\cvline{2014-2016}{\textbf{Co-founder and co-organizer, ChE Graduate Simulation Seminar Series}}
\cvlistitem{Conceptualized and started a seminar series with colleagues to address the unfulfilled need for a common collaborative and social platform for theory and computation researchers on campus.}
\cvlistitem{Developed extensive public outreach and attracted funding from ChE, Materials and Computer Science faculty, with participation from over 10 STEM fields and special attendances and keynotes from industry and national labs}

\cvline{2014-2015}{\textbf{President, India Association of Santa Barbara}}
\cvlistitem{Responsible for planning large-scale inclusive liberal arts and music festivals}
\cvlistitem{Increased participation from grad students and nearby Indian communities by 50 \%}
\end{document}